\documentclass[12pt]{article}
\usepackage[T1]{fontenc}
\usepackage[utf8]{inputenc}
\usepackage{lmodern}
\usepackage{textcomp}
\usepackage{enumitem}
\usepackage{geometry}
\geometry{margin=1in}
\begin{document}
\begin{center}
Answer Key - Computer Science - Graduate Level Assessment 22 \
\small (Answers are ordered exactly as in the question paper.)
\end{center}

\section*{Multiple Choice}
\begin{enumerate}[label=\arabic*.]
\item \textbf{Answer:} C
\\ \textit{Options:} A) Fast hardware exists to move blocks of pixels efficiently.; B) Realistic lighting and shading can be done.; C) All line segments can be displayed as straight.; D) Polygons can be filled with solid colors and textures.
\\ \textit{Note:} The statement is correct because bitmap graphics represent images as a grid of individual pixels, making it impossible to represent curved lines or diagonal segments without them appearing jagged or pixelated, thus all line segments cannot be displayed as perfectly straight.
\item \textbf{Answer:} A
\\ \textit{Options:} A) Reference counting is well suited for reclaiming cyclic structures.; B) Reference counting incurs additional space overhead for each memory cell.; C) Reference counting is an alternative to mark-and-sweep garbage collection.; D) Reference counting need not keep track of which cells point to other cells.
\\ \textit{Note:} Reference counting cannot effectively reclaim cyclic structures because the presence of mutual references between objects prevents their reference counts from ever reaching zero, leading to memory leaks.
\item \textbf{Answer:} D
\\ \textit{Options:} A) Maltego; B) BurpSuit; C) Nessus; D) Wireshark
\\ \textit{Note:} Wireshark is a widely-used network protocol analyzer that allows users to capture and examine the data packets flowing through a network, making it a powerful tool for wireless traffic sniffing.
\item \textbf{Answer:} D
\\ \textit{Options:} A) Condition codes set by every instruction; B) Variable-length encoding of instructions; C) Instructions requiring widely varying numbers of cycles to execute; D) Several different classes (sets) of registers
\\ \textit{Note:} The presence of several different classes of registers can enhance instruction-level parallelism and reduce dependency hazards, thereby facilitating rather than obstructing aggressive pipelining in an integer unit.
\item \textbf{Answer:} C
\\ \textit{Options:} A) Replaced; B) Overview; C) Changed; D) Violated
\\ \textit{Note:} A hash function produces a unique fixed-size output (hash value) for a given input, allowing any alterations to the original message to result in a different hash, thereby ensuring the integrity of the message by detecting changes.
\end{enumerate}

\section*{True / False}
\begin{enumerate}[label=\arabic*.]
\setcounter{enumi}{5}
\item \textbf{Answer:} True
\\ \textit{Options:} A) True; B) False
\\ \textit{Note:} The statement is true because network, vulnerability, and port scanning are critical techniques used in security assessments to identify potential security weaknesses and vulnerabilities in systems and networks.
\item \textbf{Answer:} False
\\ \textit{Options:} A) True; B) False
\\ \textit{Note:} The statement is False because the maximum degree of the minimal-degree interpolating polynomial for n + 1 distinct points is n, not n + 1, as a polynomial of degree n can uniquely interpolate n + 1 points if they are not all on the same polynomial of degree less than n.
\item \textbf{Answer:} False
\\ \textit{Options:} A) True; B) False
\\ \textit{Note:} The statement is false because the number of distinct functions from a set A with m elements to a set B with n elements is actually $n^m$, not n!/(n - m)!, as each element in A can be independently mapped to any of the n elements in B.
\item \textbf{Answer:} False
\\ \textit{Options:} A) True; B) False
\\ \textit{Note:} A digital signature relies on asymmetric cryptography, which requires a public-private key pair, making it incompatible with a shared-key system that uses symmetric cryptography.
\item \textbf{Answer:} True
\\ \textit{Options:} A) True; B) False
\\ \textit{Note:} Troya is considered a remote Trojan because it is designed to allow attackers to gain unauthorized access to a user's system remotely, enabling them to control and manipulate the infected device.
\end{enumerate}

\section*{Long Form Response}
\begin{enumerate}[label=\arabic*.]
\setcounter{enumi}{10}
\item \textbf{Answer:} def specified\_element(nums, N): | return result
\\ \textit{Note:} correct because it defines a function named `specified\_element` that is intended to iterate over a two-dimensional list and return a list of elements that match the specified index `N`, thereby fulfilling the requirement to extract specified elements from the input data structure.
\item \textbf{Answer:} def count\_Set\_Bits(n) : | return cnt;
\\ \textit{Note:} correct as it identifies the need for a function that iteratively or recursively calculates the total number of set bits (1 s) in the binary representations of all integers from 1 to n, although it lacks the complete implementation needed to perform the count.
\end{enumerate}

\vfill
\noindent\textit{End of Answer Key}
\end{document}